\section{Verwandte Arbeiten}
% Vorsicht vor "predatory journals" --> Diese zählen nicht als wissenschaftliche Arbeiten (Hier nachschlagen: https://dblp.org/ oder https://dl.acm.org/)
% [x] Führen sie eine kurze Literatursuche in der wissenschaftlichen Literatur zu Informationsvisualisierung und Visual Analytics nach ähnlichen Anwendungen durch.
% [ ] Diskutieren sie mindestens zwei Artikel.
% [ ] Kurze Zusammenfassung verwandter/konkurrierender Arbeiten, die das Problem auf ähnliche Weise lösen 
% [ ] Stellen sie Gemeinsamkeiten und Unterschiede dar.

% Erste Gruppe von Artikeln
% Zweite Gruppe von Artikeln
Nicht nur während der Olympischen Sommerspiele, sondern auch danach interessieren sich Menschen auf der ganzen Welt für die Ereignisse und Ergebnisse. Auch in der Wissenschaft werden die Olympischen Spiele immer wieder thematisiert und analysiert. Ein großer Teil dieser Arbeiten beschäftigt sich mit der Vorhersage neuer Medaillengewinner aufgrund historischer Entwicklungen in einer Sportart. Dabei werden vor allem Methoden der Statistik verwendet. Im Bereich der Informationsvisualisierung gibt es bisher nur wenige  Arbeiten, die eine wissenschaftliche Analyse auf den Daten der Olympischen Sommerspiele durchführen. Auf Kaggle existieren zwar einige Notebooks zu dem Thema, diese untersuchen aber keine wissenschaftliche Fragestellung, sondern stellen willkürliche Diagramme bereit. Trotzdem konnten wir zwei interessante Visual-Analytics-Anwendung finden, die ähnlich zu unserer zentralen Fragestellung, die Bewertungsgrundlage des Medaillenspiegels hinterfragen.

%Zusammenfassung, Gemeinsamkeiten & Unterschiede
\paragraph{Erster Artikel}
In \cite{relWork1} wird festgestellt, dass bestimmte Länder und Kontinente seit jeher einzelne Sportarten dominieren. In der Analyse werden zum Beispiel die Sportarten Tischtennis und Badminton genannt, welche über viele Jahren von China dominiert werden. Außerdem wird der Heimvorteil des Gastgebenden Landes, im Bezug auf den Medaillenspiegel, untersucht. In der Analyse wird eine geringe Steigerung der Anzahl gewonnener Medaillen festgestellt, welche mit der automatischen Qualifikation des Gastgebenden Landes begründet wird.

Diese Visual-Analytics-Anwendung in \cite{relWork1} soll den Anwendern verdeutlichen, dass selbst in einem solchen globalen Sportevent Dominanz und Vorteile existieren. Die Anwender sollen eine neue Perspektive auf die Ergebnisse der Olympischen Spiele erhalten, wobei ihnen die bereitgestellten Visualisierungstechniken unterstützen sollen. Die Daten in \cite{relWork1} werden von einer ähnlichen Quelle geladen, wie die historischen Daten der Olympischen Sommerspiele, welche in dieser Arbeit genutzt werden. Zur Visualisierung der Dominanz eines Landes/Kontinents in einer Sportart wird ein SunBurst-Diagramm verwendet. Diese Entscheidung wird mit der kompakten Darstellungsform begründet, da in dem SunBurst-Diagramm verschiedene Sportarten gleichzeitig dargestellt werden können. Trotzdem werden die verschiedenen Sportarten nicht alle in dem gleichen Diagramm dargestellt, sondern auf fünf Diagramme verteilt. Die Farben repräsentieren hierbei das Land/den Kontinent und die verschiedenen Jahre werden auf die Hierarchie abgebildet. Dieser Unterschied zu dieser Arbeit resultiert aus den verschiedenen Anwendungsaufgaben, welche die Visualisierung unterstützen soll. Zur Darstellung der historischen Medaillenspiegel wird in \cite{relWork1} ein Linien-Diagramm verwendet. In dieser Arbeit wurde dafür eine HeatMap verwendet, da die Zeitreihen verschiedener Länder übersichtlicher dargestellt werden kann. Die Visualisierung in einem Linien-Diagramm mit 200 Linien ist sehr unübersichtlich und konkrete Werte können nicht identifiziert werden. Dieses Problem wird in \cite{relWork1} dadurch gelöst, dass einzelne Länder hervorgehoben werden können. In der HeatMap ist dies von Anfang an möglich. Als dritte Anwendungsaufgabe in \cite{relWork1} wird ähnlich zu dieser Arbeit ein alternatives Länderranking berechnet. Dabei werden die Anzahl gewonnener Medaillen im Verhältnis zu der Anzahl der Athleten gestellt. Als Visualisierungstechnik wird ein ScatterPlot gewählt. Die Bedeutung der Kodierte Position ist bei dieser Visualisierungstechnik nicht intuitiv verständlich. In dieser Arbeit wird daher das mentale Modell des Rankings genutzt und an den Koordinatenachsen eines Parallelen Koordinatensystems abgebildet. Andererseits können in dem ScatterPlot aber sehr einfach \enquote{Ausreißer} identifiziert werden.

%Zusammenfassung, Gemeinsamkeiten & Unterschiede
\paragraph{Zweiter Artikel}
Auch in \cite{relWork2} werden die vergangenen Olympischen Sommerspiele analysiert. Die Anwender sollen dabei Einblicke in die Entwicklungen und Trends erhalten,. Dazu werden einige Fragen formuliert, welche mithilfe unterschiedlicher Visualisierungstechniken beantwortet werden sollen. Als Datenquelle wird der gleiche Datensatz wie in \cite{relWork1} verwendet.

In der ersten Anwendungsaufgabe sollen die Anwender die erfolgreichsten Länder anhand der Anzahl aller jemals gewonnenen Medaillen identifizieren. Dazu wird eine Weltkarte verwendet, in der die verschiedenen Länder mit unterschiedlichen Farben dargestellt werden. Die Farben entsprechen einem Ranking, welches nicht genauer erklärt wird. In dieser Arbeit wird der \enquote{All-Time} Medaillenspiegel nicht dargestellt, sondern nur der Medaillenspiegel der Olympischen Sommerspiele 2024 als Tabelle. Während in der Tabelle die konkrete Platzierung eines Landes leicht abgelesen werden kann, ermöglicht die Darstellung der Weltkarte globale Strukturen zu identifizieren. Beispielsweise sind auf dem Kontinent Afrika viele graue Flächen, welche Länder ohne Olympische Medaillen darstellen. Diese Art der Visualisierung bietet völlig neue Analysemöglichkeiten. In einer anderen Anwendungsaufgabe in \cite{relWork2} werden die historischen Medaillenspiegel der Länder verglichen. Dazu wird ein Linien-Diagramm verwendet, welches aber nur die fünf besten Länder darstellt. Aufgrund dieser Limitierung ist keine adäquate Analyse möglich. In dieser Arbeit wurde für diese Anwendungsaufgabe eine HeatMap verwendet, welche nicht nur die besten fünf Länder, sondern alle Länder darstellen kann. Trendanalysen sind sowohl im Linien-Diagramm durch Analyse des Verlaufs der Linien, als auch in der HeatMap durch unterschiedliche Farben, möglich. Die anderen Anwendungsaufgaben in \cite{relWork2} sind mit denen dieser Arbeit nicht zu vergleichen, weshalb darauf nicht weiter eingegangen wird.
